% Supplementary Note 3: Study 3's stated time, the gap at each time and the cell means. Numbers through macros only.
\suppnote{Study 3: stated time, the gap at each time and the cell means}

This note gives the Study~3 tests and estimates behind the main text's Results section on stated hours.

Study~3's \SThreeN{} new participants each rated nine of the ten Study~1 images for creativity and for effort on 1-to-5 scales. A label told each participant that the artist, Alex, had made the images either with design and sketching tools (the Human label; \SThreeNHuman{} participants) or by prompting an AI system (the AI label; \SThreeNAI{} participants). Each image came with a stated time that Alex had spent on it, under 1 hour, 10 hours or 100 hours, three images at each time. The label varied between participants and the stated time within them. Under the Human label Alex spent those hours drawing; under the AI label he spent them revising prompts, so under the AI label more time also meant more control over the image, and the two cannot be told apart. Study~3's registration calls the Human-label mean minus the AI-label mean the gap; a positive gap means that the AI label lowered the rating. Estimates come from the registered linear mixed model, a regression that gives each participant and each image its own baseline, fitted to each rating separately. Study~3 was registered on 8 October 2026, after its data had been collected in March 2025, after the authors had analyzed them in April 2025, and after a talk in September 2025 had reported effects of label and time (Supplementary Table~\ref{tab:supp-verdicts}).

Six $F$ tests were registered as confirmatory: label, stated time and their interaction, each on creativity and on effort. They were judged together after Holm's correction,\cite{holm1979simple} which raises each $p$-value in a set of tests so that the chance of calling any result in the set significant by mistake stays at .05. A $p$-value is the chance of a result at least this large if the factor had no effect. Two registered secondary quantities are reported as estimates, not decisions: the contrast between 100 hours and under 1 hour, and the gap at each time.

\paragraph*{Label and stated time.}
The four $F$ tests for label and for time, each on creativity and on effort, ask whether a factor changes the mean rating, and all four survived Holm's correction (registered confirmatory; supported). The label changed creativity ratings ($F$(\SThreeCreatSourceDfNum,~\SThreeCreatSourceDfDen)~=~\SThreeCreatSourceF, Holm $p$~\SThreeCreatSourcePHolm) and effort ratings ($F$(\SThreeEffSourceDfNum,~\SThreeEffSourceDfDen)~=~\SThreeEffSourceF, Holm $p$~\SThreeEffSourcePHolm). Stated time changed creativity ratings ($F$(\SThreeCreatTimeDfNum,~\SThreeCreatTimeDfDen)~=~\SThreeCreatTimeF, Holm $p$~\SThreeCreatTimePHolm) and effort ratings ($F$(\SThreeEffTimeDfNum,~\SThreeEffTimeDfDen)~=~\SThreeEffTimeF, Holm $p$~\SThreeEffTimePHolm). An $F$ test gives no direction. The gaps in Supplementary Table~\ref{tab:supp-s3cells} show the AI label lower at every time. Averaged over labels, the mean rating at 100 hours minus the mean at under 1 hour was \SThreeCreatTimeContrastEst{} points on creativity (95\% confidence interval, CI, the range of differences the data do not rule out at the 95\% level: \SThreeCreatTimeContrastLo{} to \SThreeCreatTimeContrastHi) and \SThreeEffTimeContrastEst{} on effort (\SThreeEffTimeContrastLo{} to \SThreeEffTimeContrastHi; registered secondary). Most of the creativity rise came at 100 hours: under both labels, the predicted creativity mean at 10 hours sat at about its level at under 1 hour. The predicted effort means rose at each step under both labels (Supplementary Table~\ref{tab:supp-s3cells}). No test of the rise within each label, or of equal rises under the two labels, was registered.

\begin{table}[htbp]
\centering
\caption{\textbf{Study~3: predicted mean ratings and the gap at each stated time.} Means are the registered linear mixed model's predicted means for each combination of label and stated time, on 1-to-5 scales (descriptive). The gap is the Human-label mean minus the AI-label mean, with its 95\% confidence interval (registered description: estimates the registration names, with no test of their own). Gaps are computed before rounding, so a gap can differ in its last digit from the difference of the rounded means. Each participant rated three images at each time; Human label $N$~=~\SThreeNHuman, AI label $N$~=~\SThreeNAI.}\label{tab:supp-s3cells}
\begin{tabular}{@{}llccc@{}}
\toprule
Rating & & Under 1 hour & 10 hours & 100 hours\\
\midrule
Creativity & Human label & \SThreeCreatMeanHumanUnderOneHour & \SThreeCreatMeanHumanTenHours & \SThreeCreatMeanHumanHundredHours\\
 & AI label & \SThreeCreatMeanAIUnderOneHour & \SThreeCreatMeanAITenHours & \SThreeCreatMeanAIHundredHours\\
 & Gap & \SThreeCreatGapUnderOneHourEst{} (\SThreeCreatGapUnderOneHourLo{} to \SThreeCreatGapUnderOneHourHi) & \SThreeCreatGapTenHoursEst{} (\SThreeCreatGapTenHoursLo{} to \SThreeCreatGapTenHoursHi) & \SThreeCreatGapHundredHoursEst{} (\SThreeCreatGapHundredHoursLo{} to \SThreeCreatGapHundredHoursHi)\\
\addlinespace[3pt]
Effort & Human label & \SThreeEffMeanHumanUnderOneHour & \SThreeEffMeanHumanTenHours & \SThreeEffMeanHumanHundredHours\\
 & AI label & \SThreeEffMeanAIUnderOneHour & \SThreeEffMeanAITenHours & \SThreeEffMeanAIHundredHours\\
 & Gap & \SThreeEffGapUnderOneHourEst{} (\SThreeEffGapUnderOneHourLo{} to \SThreeEffGapUnderOneHourHi) & \SThreeEffGapTenHoursEst{} (\SThreeEffGapTenHoursLo{} to \SThreeEffGapTenHoursHi) & \SThreeEffGapHundredHoursEst{} (\SThreeEffGapHundredHoursLo{} to \SThreeEffGapHundredHoursHi)\\
\bottomrule
\end{tabular}
\end{table}

\paragraph*{Whether the gap changed with time.}
The label-by-time interaction asks whether the gap changed with time. On neither rating did it survive Holm's correction (registered confirmatory): on creativity, $F$(\SThreeCreatSourceByTimeDfNum,~\SThreeCreatSourceByTimeDfDen)~=~\SThreeCreatSourceByTimeF{} and Holm $p$~\SThreeCreatSourceByTimePHolm; on effort, $F$(\SThreeEffSourceByTimeDfNum,~\SThreeEffSourceByTimeDfDen)~=~\SThreeEffSourceByTimeF{} and Holm $p$~\SThreeEffSourceByTimePHolm. In the ordinal check, a registered secondary analysis that treats the five answers as ordered categories rather than evenly spaced numbers, the effort interaction came nearer the threshold (uncorrected $p$~\SThreeEffSourceByTimeOrdP, Holm $p$~\SThreeEffSourceByTimeOrdPHolm), but still fell short of significance after correction. On every term the ordinal check reached the same verdict as the linear model. The September 2025 talk had also reported an effort interaction from an earlier analysis that still included participants the registration later excluded; the registered model does not reproduce it. A nonsignificant interaction does not show that the gap stayed the same, and no test of equal gaps across times was registered.

The registered decision on both ratings, that the gap persists at 100 hours, is reported in the main text. Part of that decision, also registered confirmatory, is an equivalence test, which asks whether the 100-hour gap is shown to lie inside $\pm$\SourceBound{} points (half of Study~1's creativity penalty). The rule asked that this test be repeated in a fuller model that lets the time effect vary from person to person and the label's effect vary from image to image. It waived that requirement if the fuller model was singular, meaning that it estimated a variance or correlation at the edge of its possible range, and on both ratings it was. The fuller model's 90\% CIs for the 100-hour gap (\SThreeCreatGapHundredSlopeLoNinety{} to \SThreeCreatGapHundredSlopeHiNinety{} on creativity, \SThreeEffGapHundredSlopeLoNinety{} to \SThreeEffGapHundredSlopeHiNinety{} on effort) also reach past the bound, so the waiver changes no verdict.

\paragraph*{The gap at each time.}
As registered description, the creativity gap was \SThreeCreatGapUnderOneHourEst, \SThreeCreatGapTenHoursEst{} and \SThreeCreatGapHundredHoursEst{} points at under 1, 10 and 100 hours, and the effort gap was \SThreeEffGapUnderOneHourEst, \SThreeEffGapTenHoursEst{} and \SThreeEffGapHundredHoursEst{} (Supplementary Table~\ref{tab:supp-s3cells}; main-text Fig.~2a plots both). Each 95\% CI lay above zero except that of the effort gap at 10 hours. By the registration, a gap significant at one time and not at another decides nothing about whether time closes the gap.

Each effort gap was smaller than Study~1's effort penalty of \SOneHOneEffDiff{} points, a comparison across studies that is untested. The two studies differ in their participants, in whether a time was stated and in the wording of the Human label, which in Study~3 named the drawing tools (iPad apps and Adobe Photoshop); any of these, or chance, could make the difference. Study~3 registered no test of effort against creativity, and with time stated its two gaps were close: at 10 hours the effort gap was the smaller. The larger drop on effort therefore rests on Study~1.

\paragraph*{Comparisons without a test.}
On effort, the AI-label mean at 100 hours (\SThreeEffMeanAIHundredHours) sat above the Human-label mean at under 1 hour (\SThreeEffMeanHumanUnderOneHour); on creativity it sat below (\SThreeCreatMeanAIHundredHours{} against \SThreeCreatMeanHumanUnderOneHour). These two cells differ in both label and stated time, and no test compares them, so the comparison is descriptive.
